\documentclass[12pt,a4paper]{article}
\usepackage[UTF8]{ctex}
\usepackage{amsmath,amssymb,amsthm}
\usepackage{geometry}

\geometry{margin=2.5cm}

\title{从关节速度控制角度理解P控制器的稳态误差}
\author{第11章补充说明}
\date{}

\begin{document}

\maketitle

\section{速度输入控制回顾}

在速度输入的运动控制中，控制律为：
\begin{equation}
\dot{\theta} = \dot{\theta}_d + K_p(\theta_d - \theta)
\end{equation}
其中：
\begin{itemize}
\item $\dot{\theta}$是实际关节速度（控制输出）
\item $\dot{\theta}_d$是期望关节速度
\item $\theta_d$是期望关节位置
\item $\theta$是实际关节位置
\item $K_p > 0$是比例增益
\item $\theta_e = \theta_d - \theta$是位置误差
\end{itemize}

\section{从速度角度分析稳态误差}

\subsection{情况1：跟踪固定位置（$\dot{\theta}_d = 0$）}

假设期望位置固定，即$\dot{\theta}_d = 0$，控制律变为：
\begin{equation}
\dot{\theta} = K_p(\theta_d - \theta) = K_p \theta_e
\end{equation}

\subsubsection{无干扰情况}

如果系统没有持续干扰（如重力），在稳态时：
\begin{equation}
\dot{\theta}_{ss} = 0
\end{equation}

因此：
\begin{equation}
K_p \theta_{e,ss} = 0
\end{equation}

由于$K_p > 0$，必须有：
\begin{equation}
\theta_{e,ss} = 0
\end{equation}

\textbf{结论：}在无干扰情况下，速度输入的P控制器可以达到零稳态误差。

\subsubsection{有持续干扰情况（重力）}

现在考虑有重力的情况。虽然速度控制器直接控制速度，但实际系统仍然受到重力的影响。

\textbf{关键问题：}速度控制器假设可以精确控制速度，但实际物理系统有动力学：
\begin{equation}
M\ddot{\theta} = \tau - mgr\cos\theta
\end{equation}

如果速度控制器试图维持$\dot{\theta} = 0$，但存在重力$mgr\cos\theta$，则：
\begin{itemize}
\item 要保持$\dot{\theta} = 0$，需要$\tau = mgr\cos\theta$
\item 但速度控制器只根据误差$\theta_e$产生速度指令
\item 如果$\theta_e = 0$，则$\dot{\theta} = 0$，但实际系统在重力作用下会加速
\item 因此，系统需要非零误差来产生非零速度，以"对抗"重力
\end{itemize}

\textbf{更准确的分析：}

实际系统动力学：
\begin{equation}
M\ddot{\theta} = \tau - mgr\cos\theta
\end{equation}

速度控制器产生速度指令：
\begin{equation}
\dot{\theta}_{\text{cmd}} = K_p(\theta_d - \theta)
\end{equation}

但实际速度$\dot{\theta}$由动力学决定。在稳态时，如果$\dot{\theta} = 0$且$\ddot{\theta} = 0$：
\begin{equation}
0 = \tau_{ss} - mgr\cos\theta_{ss}
\end{equation}

如果速度控制器试图维持$\dot{\theta} = 0$，但$\theta_e = 0$，则$\dot{\theta}_{\text{cmd}} = 0$，无法产生抵消重力的力矩。

\textbf{因此，系统会在非零误差处达到平衡：}
\begin{equation}
\dot{\theta}_{\text{cmd}} = K_p \theta_{e,ss} \neq 0
\end{equation}

这个非零速度指令产生的力矩（通过执行器）正好抵消重力。

\subsection{情况2：跟踪运动轨迹（$\dot{\theta}_d \neq 0$）}

假设期望位置在变化，$\dot{\theta}_d \neq 0$。

控制律：
\begin{equation}
\dot{\theta} = \dot{\theta}_d + K_p(\theta_d - \theta)
\end{equation}

\subsubsection{理想情况：无干扰}

如果系统能精确跟踪，在稳态时：
\begin{equation}
\dot{\theta} = \dot{\theta}_d, \quad \theta = \theta_d
\end{equation}

因此：
\begin{equation}
\theta_{e,ss} = 0
\end{equation}

\textbf{结论：}在理想无干扰情况下，可以跟踪运动轨迹且无稳态误差。

\subsubsection{实际情况：有干扰}

如果存在持续干扰（如重力），系统无法完全跟踪期望轨迹。

假设期望轨迹是匀速运动：$\dot{\theta}_d = \text{常数}$，$\theta_d(t) = \theta_{d,0} + \dot{\theta}_d t$。

在稳态时，如果实际速度$\dot{\theta} = \dot{\theta}_d$，但存在位置误差$\theta_e \neq 0$：
\begin{equation}
\dot{\theta} = \dot{\theta}_d + K_p \theta_e = \dot{\theta}_d
\end{equation}

这要求$K_p \theta_e = 0$，即$\theta_e = 0$。

但如果存在重力干扰，要保持$\dot{\theta} = \dot{\theta}_d$，需要额外的力矩来抵消重力。如果$\theta_e = 0$，速度控制器无法提供这个力矩。

\textbf{因此，系统会存在稳态误差：}
\begin{equation}
\theta_{e,ss} \neq 0
\end{equation}

这个误差产生的速度修正项$K_p \theta_{e,ss}$通过执行器转换为力矩，抵消重力。

\section{具体数值示例}

\subsection{示例1：固定位置跟踪（有重力）}

\textbf{系统参数：}
\begin{itemize}
\item 期望位置：$\theta_d = 0$ rad
\item 初始位置：$\theta(0) = 0.2$ rad
\item 比例增益：$K_p = 10$ rad/s/rad
\item 重力力矩：$mgr = 5$ N·m（假设在$\theta = 0$处）
\item 执行器特性：速度指令转换为力矩，转换系数$\alpha = 2$ N·m/(rad/s)
\end{itemize}

\textbf{控制律：}
\begin{equation}
\dot{\theta}_{\text{cmd}} = 10(0 - \theta) = -10\theta
\end{equation}

\textbf{稳态分析：}

在稳态时，如果$\dot{\theta} = 0$，系统动力学要求：
\begin{equation}
\tau_{ss} = mgr = 5 \text{ N·m}
\end{equation}

速度控制器产生的速度指令为：
\begin{equation}
\dot{\theta}_{\text{cmd}} = -10 \theta_{ss}
\end{equation}

这个速度指令通过执行器转换为力矩：
\begin{equation}
\tau_{ss} = \alpha \cdot \dot{\theta}_{\text{cmd}} = 2 \times (-10 \theta_{ss}) = -20 \theta_{ss}
\end{equation}

平衡条件：
\begin{equation}
-20 \theta_{ss} = 5
\end{equation}

因此：
\begin{equation}
\theta_{ss} = -\frac{5}{20} = -0.25 \text{ rad}
\end{equation}

稳态误差：
\begin{equation}
\theta_{e,ss} = 0 - (-0.25) = 0.25 \text{ rad} \approx 14.3^\circ
\end{equation}

\textbf{验证：}
\begin{align}
\dot{\theta}_{\text{cmd}} &= -10 \times (-0.25) = 2.5 \text{ rad/s} \\
\tau_{ss} &= 2 \times 2.5 = 5 \text{ N·m} \quad \checkmark
\end{align}

\textbf{结论：}系统在$\theta = -0.25$ rad处达到平衡，稳态误差为$0.25$ rad。

\subsection{示例2：增大增益的影响}

如果增大$K_p = 50$ rad/s/rad，其他参数不变：

\begin{equation}
\tau_{ss} = 2 \times (-50 \theta_{ss}) = -100 \theta_{ss} = 5
\end{equation}

因此：
\begin{equation}
\theta_{ss} = -\frac{5}{100} = -0.05 \text{ rad}
\end{equation}

稳态误差：
\begin{equation}
\theta_{e,ss} = 0.05 \text{ rad} \approx 2.9^\circ
\end{equation}

\textbf{结论：}增大$K_p$可以减小稳态误差（从$0.25$ rad减小到$0.05$ rad），但不能完全消除。

\subsection{示例3：跟踪运动轨迹}

\textbf{系统参数：}
\begin{itemize}
\item 期望速度：$\dot{\theta}_d = 1$ rad/s（匀速运动）
\item 期望位置：$\theta_d(t) = t$ rad
\item 比例增益：$K_p = 10$ rad/s/rad
\item 重力力矩：$mgr = 5$ N·m
\item 执行器转换系数：$\alpha = 2$ N·m/(rad/s)
\end{itemize}

\textbf{控制律：}
\begin{equation}
\dot{\theta}_{\text{cmd}} = 1 + 10(\theta_d - \theta)
\end{equation}

\textbf{稳态分析：}

假设在稳态时，实际速度$\dot{\theta} = 1$ rad/s（跟踪期望速度），但存在位置误差$\theta_e$。

速度指令：
\begin{equation}
\dot{\theta}_{\text{cmd}} = 1 + 10 \theta_e
\end{equation}

实际速度由动力学决定。如果$\dot{\theta} = 1$，则$\ddot{\theta} = 0$，因此：
\begin{equation}
\tau_{ss} = mgr = 5 \text{ N·m}
\end{equation}

速度指令产生的力矩：
\begin{equation}
\tau_{ss} = \alpha \cdot \dot{\theta}_{\text{cmd}} = 2(1 + 10 \theta_e) = 2 + 20 \theta_e
\end{equation}

平衡条件：
\begin{equation}
2 + 20 \theta_e = 5
\end{equation}

因此：
\begin{equation}
\theta_e = \frac{5 - 2}{20} = 0.15 \text{ rad} \approx 8.6^\circ
\end{equation}

\textbf{结论：}即使在跟踪运动轨迹时，如果存在持续干扰，P控制器也会产生稳态误差。

\section{从速度角度的直观理解}

\subsection{关键洞察}

在速度输入控制中，P控制器的稳态误差可以这样理解：

\begin{enumerate}
\item \textbf{速度指令的作用：}控制律$\dot{\theta} = \dot{\theta}_d + K_p \theta_e$产生速度指令
\item \textbf{速度到力矩的转换：}速度指令通过执行器转换为力矩
\item \textbf{力矩平衡：}在稳态时，这个力矩必须抵消持续干扰（如重力）
\item \textbf{误差的必要性：}如果$\theta_e = 0$，则速度修正项$K_p \theta_e = 0$，无法产生抵消干扰的力矩
\item \textbf{因此：}必须存在非零误差$\theta_e \neq 0$来产生非零速度修正，进而产生抵消干扰的力矩
\end{enumerate}

\subsection{类比：水流系统}

想象一个水流系统：
\begin{itemize}
\item 期望水位：$h_d$
\item 实际水位：$h$
\item 控制律：流量$Q = K_p(h_d - h)$
\item 持续泄漏：$Q_{\text{leak}} = \text{常数}$
\end{itemize}

在稳态时，要保持水位不变，需要：
\begin{equation}
Q = Q_{\text{leak}}
\end{equation}

如果$h = h_d$，则$Q = 0$，无法抵消泄漏，水位会下降。

因此，必须存在水位误差$h_d - h \neq 0$，使得$Q = K_p(h_d - h) = Q_{\text{leak}}$。

\section{为什么速度控制中也可能有稳态误差？}

虽然速度控制器直接控制速度，但：

\begin{enumerate}
\item \textbf{物理系统有动力学：}实际系统不是理想的速度源，而是有质量、惯性的动力学系统
\item \textbf{持续干扰存在：}重力、摩擦力等持续干扰仍然作用在系统上
\item \textbf{速度到力矩的转换：}速度指令需要通过执行器转换为力矩，才能影响系统
\item \textbf{力矩平衡要求：}在稳态时，产生的力矩必须抵消干扰
\item \textbf{误差的必要性：}如果误差为零，速度修正项为零，无法产生抵消干扰的力矩
\end{enumerate}

\section{总结}

\textbf{从关节速度控制角度理解稳态误差：}

\begin{enumerate}
\item \textbf{控制律：}$\dot{\theta} = \dot{\theta}_d + K_p(\theta_d - \theta)$

\item \textbf{稳态误差的原因：}
   \begin{itemize}
   \item 存在持续干扰（如重力）
   \item 速度修正项$K_p \theta_e$通过执行器转换为力矩
   \item 在稳态时，这个力矩必须抵消干扰
   \item 如果$\theta_e = 0$，则$K_p \theta_e = 0$，无法产生抵消干扰的力矩
   \item 因此，必须存在$\theta_e \neq 0$
   \end{itemize}

\item \textbf{稳态误差表达式：}
   \begin{equation}
   \theta_{e,ss} \approx \frac{\text{持续干扰力矩}}{\alpha \cdot K_p}
   \end{equation}
   其中$\alpha$是速度到力矩的转换系数。

\item \textbf{减小稳态误差的方法：}
   \begin{itemize}
   \item 增大$K_p$（但不能无限增大，受物理限制）
   \item 使用PI控制器（加入积分项）
   \item 前馈补偿干扰（如果干扰已知）
   \end{itemize}
\end{enumerate}

\textbf{关键点：}即使在速度输入控制中，如果存在持续干扰，P控制器也会产生稳态误差，因为需要非零误差来产生抵消干扰的力矩。

\end{document}


